\subsection{Erweiterung durch Zusatzfunktionen}
Neben der reinen Zeitanzeige wird das System um zusätzliche informative Funktionen erweitert. 
Hierzu zählen die Anzeige von Wetter- und Temperaturdaten, die Tageszeit, die aktuelle Mondphase sowie die Integration eines Innenraumsensors zur Erfassung von Temperatur und Luftfeuchtigkeit. 
Die Daten werden in definierten Zeitintervallen aktualisiert und anschließend in abstrahierter Form dargestellt. 
Der Ablauf ist in Abbildung~\ref{fig:erweiterung} dargestellt.

\begin{figure}[H]
\centering
\begin{tikzpicture}[
    node distance=1.0cm,
    startstop/.style={rectangle, rounded corners, minimum width=4cm, draw=black},
    process/.style={rectangle, minimum width=4.8cm, draw=black},
    decision/.style={diamond, aspect=2.3, draw=black},
    io/.style={trapezium, trapezium left angle=70, trapezium right angle=110, minimum width=4.6cm, minimum height=0.8cm, text centered, draw=black},
    arrow/.style={->, thick}
]

\node (start) [startstop] {Zusatzfunktionen};
\node (weather) [process, below=of start] {Open-Meteo-API Daten abfragen};
\node (sensor) [process, below=of weather] {Sensorwerte abfragen};
\node (display) [io, below=of sensor] {Zusatzfunktionen anzeigen};
\node (end) [startstop, below=of display] {Ende};

\draw [arrow] (start) -- (weather);
\draw [arrow] (weather) -- (sensor);
\draw [arrow] (sensor) -- (display);
\draw [arrow] (display) -- (end);

\end{tikzpicture}
\caption{Schematischer Programmablaufplan zur Erzeugung der Anzeige für die WordClock, welche die informativen Zusatzfunktionen darstellt.}
\label{fig:erweiterung}
\end{figure}